\documentclass[11pt]{article}

\usepackage[margin=1in]{geometry}
\usepackage{amsmath, amssymb, amsfonts}
\usepackage{graphicx}
\usepackage{booktabs}
\usepackage[colorlinks=true, linkcolor=blue, citecolor=blue, urlcolor=blue]{hyperref}
\usepackage[round]{natbib}
\usepackage{xcolor}

%TODO marker for pending campaign results
\newcommand{\TBD}[1]{\textcolor{red}{[TBD: #1]}}

\title{NeuroFoil: An Open Pipeline Coupling Real-OpenFOAM RANS Data with\\
Calibrated Neural Surrogates for Two-Dimensional Airfoil Design}

\author{%
  BLACKNIGHT (cloud12)\thanks{Independent researcher. Code and data:
  \url{https://github.com/BLACKNIGHT-cloud12/NeuroFoil-Core}}%
}
\date{October 2026}

\begin{document}
\maketitle

\begin{abstract}
Neural surrogates can collapse the cost of aerodynamic design exploration from
minutes-per-evaluation (RANS CFD) to sub-millisecond inference, but published
surrogate pipelines frequently entangle three failure modes: training data of
uncertain provenance, train/test contamination across airfoil geometries, and
uncertainty estimates that are neither calibrated nor evaluated. We present
NeuroFoil, an open end-to-end pipeline that addresses all three. A fully
automated OpenFOAM workflow (steady incompressible k--$\omega$ SST on
parametric six-block structured C-grids, $y^+ \approx 1$, potential-flow
initialization) generates a campaign dataset of \TBD{N} simulation records
spanning \TBD{21} airfoils, four Reynolds numbers ($1$--$6\times10^6$), and six
angles of attack ($-6^\circ$ to $14^\circ$), with every row carrying its honest
convergence status. Residual MLP and surface-pressure field-decoder networks are
trained with \emph{airfoil-grouped} splits so that reported test metrics measure
extrapolation to unseen geometries rather than memorization, alongside
random-row interpolation baselines. Epistemic uncertainty is quantified with
split-conformal calibration whose empirical coverage is measured on held-out
airfoils, and out-of-distribution queries are flagged by a Mahalanobis guard on
the training manifold. The trained networks deploy through a pure-NumPy runtime
($<0.5$\,ms per query, no deep-learning framework required) with a one-click
OpenFOAM verification loop. Solver settings are verified against classical
NACA~0012 benchmarks: the computed lift slope ($6.29$\,rad$^{-1}$) matches
thin-airfoil theory and experimental measurements, and a three-level grid study
quantifies the production-mesh discretization band.
\end{abstract}

\section{Introduction}
\label{sec:intro}

Surrogate-assisted design promises orders-of-magnitude speedups over
direct CFD evaluation, but the value of a surrogate rests entirely on the
credibility of its training data and the honesty of its evaluation. Three
failure modes recur in surrogate-aerodynamics pipelines:
\begin{enumerate}
  \item \textbf{Data provenance.} Datasets presented as ``CFD-derived'' are
  sometimes synthesized from textbook formulas, making downstream accuracy
  claims circular: the surrogate is evaluated against the same analytic model
  that generated its training data.
  \item \textbf{Split contamination.} Parametric airfoil datasets naturally
  repeat each geometry across many operating points; a random row-wise
  train/test split therefore places near-duplicates on both sides and
  overstates generalization.
  \item \textbf{Uncalibrated uncertainty.} Hand-tuned variance formulas are
  presented as confidence intervals without any coverage measurement, and
  out-of-distribution (OOD) detectors are deployed without detection metrics.
\end{enumerate}

NeuroFoil is built to be auditable against each of these. Its dataset is
produced by a scripted OpenFOAM campaign whose solver settings are verified
against classical benchmarks (Section~\ref{sec:verification}); its models are
trained and reported under airfoil-grouped splits with interpolation baselines
(Section~\ref{sec:setup}); its intervals are split-conformal with measured
coverage, and its OOD guard is evaluated on held-out geometries
(Section~\ref{sec:uq}). Everything --- case dictionaries, mesh generator,
campaign driver, training notebook, and a framework-free inference runtime ---
is released under a permissive license.

Our contributions are:
\begin{itemize}
  \item an open, reproducible, Windows/Linux-portable RANS data-generation
  pipeline with verified numerical settings (under-relaxation and
  potential-flow initialization are shown to be load-bearing);
  \item a leakage-free evaluation protocol for parametric airfoil surrogates
  that separates interpolation from shape extrapolation;
  \item split-conformal uncertainty for scalar coefficients and
  161-point surface-pressure fields with empirical coverage audit;
  \item a deployment path that runs the trained network in pure NumPy,
  enabling framework-free, sub-millisecond inference in engineering tools.
\end{itemize}

\section{Related Work}
\label{sec:related}

\textbf{Airfoil parameterization.} The NACA 4-digit family remains the
canonical testbed for aerodynamic methods \citep{abbott1959}, and the
Class-Shape Transformation (CST) of \citet{kulfan2008} provides a compact,
order-expandable representation that we use as surrogate input features.

\textbf{Turbulence modeling.} We adopt the shear-stress-transport
k--$\omega$ model of \citet{menter1994}, the de-facto industrial standard for
attached and mildly separated external aerodynamics. OpenFOAM
\citep{weller1998} supplies the finite-volume solver infrastructure.

\textbf{Experimental benchmarks.} The low-speed NACA~0012 measurements of
\citet{ladson1988} and the tabulations of \citet{abbott1959} anchor our
solver-verification study (Section~\ref{sec:verification}).

\textbf{Conformal prediction.} Split conformal prediction
\citep{papadopoulos2002, vovk2005} yields distribution-free interval estimates
with finite-sample marginal coverage guarantees under exchangeability, and has
recently been applied to regression surrogates; we apply it both to scalar
force coefficients and pointwise to surface-pressure fields.

\textbf{OOD detection.} We use the Mahalanobis-distance score of
\citet{lee2018} on the training manifold, evaluated against held-out
geometries rather than merely asserted.

\TBD{Final literature pass: surrogate modeling in aerodynamics
(mlp/GP/field-prediction literature), conformal applications in engineering
UQ, and positioned comparison to XFOIL-based datasets.}

\section{Methods}
\label{sec:methods}

\subsection{Geometry and parameterization}
Airfoils are generated from the NACA 4-digit equations with the closed
trailing-edge coefficient $a_4=-0.1036$ and cosine chordwise clustering
(161 points per surface). Each surface is fitted with a fifth-order Kulfan CST
Bernstein expansion \citep{kulfan2008} with class exponents
$N_1=0.5$, $N_2=1.0$; the twelve CST weights, together with $(m, p, t)$,
form the geometry features of the surrogate. A fixed cosine grid
$x_i = \tfrac12(1-\cos \beta_i)$, $\beta_i = i\pi/160$, standardizes all
surface-pressure fields.

\subsection{Automated RANS pipeline}
The pipeline (OpenFOAM v2606 \citep{weller1998}) builds a parametric six-block
2D structured C-grid per case: farfield radius 15 chords, wake extension 15
chords, span 0.1, and a wall-normal first-cell height of
$4\times10^{-6}$ chords (geometric grading solved in closed form from the
requested first-cell height), giving $y^+\approx1$ at
$\mathrm{Re}=6\times10^6$. Boundary patches: fixed-velocity inlet over the
upstream arc and lateral boundaries, fixed-pressure outlet downstream,
no-slip airfoil wall, empty front/back for the 2D constraint.

The case runs steady incompressible simpleFoam with k--$\omega$ SST
\citep{menter1994}. Two settings proved load-bearing during verification:
(i)~\emph{under-relaxation} p=0.1, U=0.3, k/$\omega$=0.5 --- the conventional
p=0.3/U=0.7 combination produces a sustained force limit cycle (lift
oscillating by $\mathcal{O}(10^2)$) on these grids; and (ii)~\emph{potential
flow initialization} via potentialFoam, which removes a large startup
transient (lift swings of $\mathcal{O}(10^3)$ with uniform initialization) and
accelerates convergence. Runs proceed to 3000 iterations with residual targets
p$<10^{-5}$, U,k,$\omega$$<10^{-6}$; convergence is accepted when the
sliding-window standard deviations of the last 50--100 force-coefficient
samples satisfy $\sigma_{C_L}<5\times10^{-4}$,
$\sigma_{C_D}<2\times10^{-4}$, $\sigma_{C_m}<5\times10^{-4}$. Every dataset
row records its true convergence status; no row is filtered silently.

Freestream turbulence follows low-turbulence wind-tunnel practice
(intensity 0.1\%, eddy-viscosity ratio 0.1). Surface-pressure fields are
extracted from the airfoil patch face centroids and resampled onto the 161-point
cosine grid; force coefficients come from the forceCoeffs function object.

\subsection{Surrogate architectures}
Two networks are trained (PyTorch \citep{paszke2019}, AdamW, cosine schedule,
all RNG seeded):
\begin{itemize}
  \item \textbf{Scalar surrogate}: a deep residual MLP (128-wide, two residual
  blocks with batch normalization and GELU, 64-wide head) mapping 17 features
  (m, p, t, 12 CST weights, $\alpha$, Re) to $(C_L, C_D, C_m)$.
  \item \textbf{Field decoder}: a plain MLP (128--256--512) mapping the same
  features to the 322-point surface-pressure vector (upper then lower
  surface); $\Delta C_p(x)$ follows by subtraction, and integrated
  $\int_0^1 \Delta C_p\,dx$ provides a physics-consistency check against the
  predicted $C_L$.
\end{itemize}

\subsection{Calibrated uncertainty and OOD detection}
\label{sec:uq}
\textbf{Split conformal.} The validation airfoils (unseen in training) form a
calibration set; the conformal quantile is the
$\lceil (n+1)(1-\alpha)\rceil/n$ quantile of absolute residuals, applied to
each scalar target and pointwise across the chord for both Cp surfaces. The
training notebook reports empirical coverage on untouched test airfoils.
\textbf{Mahalanobis guard.} Following \citet{lee2018}, the score
$D_M = \sqrt{(\mathbf{x}-\boldsymbol{\mu})^T \boldsymbol{\Sigma}^{-1}
(\mathbf{x}-\boldsymbol{\mu})}$ is computed on the training manifold
(ridged covariance); envelope violations in $(m,p,t,\alpha,\mathrm{Re})$ are
flagged directly. Detection performance on held-out airfoils is reported in
Section~\ref{sec:results}.

\subsection{Multidisciplinary design optimization}
An SLSQP loop \citep{kraft1988} maximizes $L/D$ subject to a lift target and
wingbox structural constraints (front/rear spar depths at 20\%/65\% chord,
inter-spar fuel-bay area), evaluated from the surrogate in closed form; each
optimization completes in well under a second on a laptop CPU.

\section{Experimental setup}
\label{sec:setup}

\textbf{Campaign design.} A Latin hypercube over $(m, p, t)$
\citep{mckay1979}, discretized to valid NACA codes and anchored with NACA
0012/2412/4412, yields \TBD{21} airfoils; crossed with
$\mathrm{Re}\in\{1,2,3,6\}\times10^6$ and
$\alpha\in\{-6^\circ,-2^\circ,2^\circ,6^\circ,10^\circ,14^\circ\}$ gives
\TBD{504} cases at the production (medium) mesh. The driver is resume-safe and
parallel (multiple WSL/native solver instances), and trims solver artifacts
after extraction.

\textbf{Evaluation protocol.} Airfoil-grouped splits hold out
$\approx$15\% of geometries each for calibration and test; a random-row split
is trained identically to quantify the interpolation--extrapolation gap.
Metrics: $R^2$, RMSE, MAE per scalar target; flattened and per-surface field
$R^2$; integrated-lift physics consistency; conformal coverage; Mahalanobis
detection rates on unseen geometries. All figures and tables are regenerated
by a single script from the released artifacts.

\section{Solver verification}
\label{sec:verification}

Before campaign production, the pipeline was verified on NACA 0012 at
$\mathrm{Re}=6\times10^6$ (18{,}048-cell medium mesh):

\begin{table}[h]
\centering
\caption{Pipeline verification, NACA 0012, Re $=6\times10^6$.}
\label{tab:verify}
\begin{tabular}{lcccc}
\toprule
Case & $C_L$ & $C_D$ & $C_{m,c/4}$ & Status \\
\midrule
$\alpha=0^\circ$ & $-0.0042$ & 0.00879 & $+0.0008$ & converged \\
$\alpha=4^\circ$ & $+0.4388$ & 0.00951 & $+0.0024$ & converged
($\sigma_{C_L}=2.2\times10^{-4}$) \\
\bottomrule
\end{tabular}
\end{table}

The $\alpha=4^\circ$ result implies a lift slope of $6.29$\,rad$^{-1}$,
consistent with thin-airfoil theory ($2\pi$) and with the measured NACA 0012
slope \citep{abbott1959, ladson1988}; the fully-turbulent drag level sits
above the transitional experimental value, as expected without a transition
model. The symmetric case returns $C_L\approx0$ and $C_{m}\approx0$ as
required. A three-level grid study quantifies the production-mesh band:

\begin{table}[h]
\centering
\caption{Grid convergence, NACA 0012, $\alpha=4^\circ$, Re $=6\times10^6$.}
\label{tab:grid}
\begin{tabular}{lccc}
\toprule
Mesh & Cells & $C_L$ & $C_D$ \\
\midrule
coarse & 8{,}640 & 0.43133 & 0.010730 \\
medium (production) & 18{,}048 & 0.43875 & 0.009509 \\
fine & 36{,}000 & 0.45908 & 0.008309 \\
\bottomrule
\end{tabular}
\end{table}

\section{Results}
\label{sec:results}

\TBD{Populate from the campaign dataset and training-notebook metrics:
(a)~grouped-split vs random-split scalar metrics;
(b)~field $R^2$ and integrated-lift consistency;
(c)~conformal coverage on unseen airfoils;
(d)~Mahalanobis detection on held-out geometries;
(e)~per-case wall-clock vs surrogate latency (cost table);
(f)~MDAO case study with one-click CFD verification of the optimum;
(g)~dataset overview (status distribution, polar sanity across Re).}

\section{Limitations}
\label{sec:limits}

The platform is 2D and incompressible (valid for $M\le0.3$); it resolves no
induced drag, laminar--turbulent transition, or transonic effects. The
linear-eddy-viscosity closure delays separation onset under strong adverse
gradients, and near-stall cases (particularly $\alpha=\pm14^\circ$) may fail
the stationarity criterion --- such rows are retained with their true status
rather than filtered. The Mahalanobis guard is calibrated to the NACA
4-digit manifold and is not a general-purpose OOD method beyond it.

\section{Conclusion}
\label{sec:conclusion}

\TBD{Final conclusions once results are populated.}

\bibliographystyle{plainnat}
\bibliography{refs}

\end{document}
